% ============================================================================
\section{Pendahuluan}
\label{sec:pendahuluan}
% ============================================================================

PPh Pasal 21 dipotong setiap bulan oleh pemberi kerja atas penghasilan pegawainya, sehingga ketepatan perhitungannya menyentuh setiap slip gaji. Aturannya juga sering berubah. Sampai tahun pajak 2023, pemotongan bulanan memakai metode disetahunkan menurut PER-16/PJ/2016~\cite{per16_2016}. Sejak 1 Januari 2024, PP 58/2023~\cite{pp58_2023} dan PMK 168/2023~\cite{pmk168_2023} menggantinya dengan tarif efektif rata-rata (TER) bulanan, ditambah perhitungan ulang setahun di masa pajak terakhir. Tahun 2025 dan 2026 menambahkan fasilitas PPh 21 ditanggung pemerintah (DTP) bagi sektor tertentu, yang salah satunya ditetapkan sesudah masa pajak yang dicakupnya dimulai~\cite{pmk10_2025,pmk72_2025,pmk105_2025}. Parameter BPJS seperti batas upah Jaminan Pensiun juga diperbarui setiap tahun. Di atas semua itu, setiap pemberi kerja memiliki kebijakan payroll sendiri, misalnya prorata kehadiran, kenaikan gaji di tengah bulan, cara menghitung THR, dan komponen tunjangan. Kebijakan itu menentukan berapa penghasilan bruto yang kemudian dikenai pajak.

Dalam praktik, aturan pemerintah dan kebijakan perusahaan ditanam bercampur di rumus \en{spreadsheet} atau kode program. Audit awal atas \en{spreadsheet} payroll sebuah perusahaan (Subbab~\ref{subsec:data}) menemukan tiga masalah. Konstanta regulasi tertanam di 254 sel rumus. Kebijakan perusahaan mengklasifikasikan lembur sebagai penghasilan tidak teratur, padahal regulasi menetapkannya sebagai penghasilan teratur. THR dipajaki pada masa yang tidak sesuai dengan saat terutangnya. Survei atas tujuh implementasi PPh 21 \en{open source}, berupa pustaka dan aplikasi, juga menemukan kesalahan terverifikasi di setiap implementasi. Contohnya lapisan tarif yang diperlakukan sebagai lebar, bukan batas kumulatif; premi BPJS pemberi kerja yang tidak masuk bruto; dan lebih bayar Desember yang dipotong menjadi nol. Temuan serupa dilaporkan pada implementasi resmi di negara lain. Formalisasi perhitungan tunjangan keluarga Prancis dengan bahasa Catala menemukan \en{bug} pada implementasi resminya~\cite{merigoux2021catala}, dan mesin pajak penghasilan Prancis sulit dipelihara karena aturannya bergantung pada bahasa dan \en{compiler} warisan~\cite{merigoux2021mlang}.

Masalah ini dapat dirumuskan dalam empat butir:
\begin{enumerate}[leftmargin=*,nosep]
  \item Setiap perubahan regulasi menuntut perubahan rumus, dan tidak ada jaminan bahwa perubahan itu tidak menggeser hasil tahun lain.
  \item Batas antara aturan wajib pemerintah dan kebijakan perusahaan tidak eksplisit, sehingga pelanggaran aturan wajib tidak terdeteksi.
  \item Angka di slip gaji tidak dapat ditelusuri ke pasal yang mendasarinya, sehingga sulit diaudit.
  \item Definisi melingkar seperti \en{gross-up}, ketika tunjangan pajak sama dengan PPh atas penghasilan yang memuat tunjangan itu sendiri, diselesaikan dengan rumus turunan manual yang tidak lagi berlaku ketika struktur tarif berubah ke TER.
\end{enumerate}
Menguji kebenaran perangkat lunak pajak pun sulit, karena spesifikasi eksplisit dan \en{test oracle} umumnya tidak tersedia~\cite{tizpaz2023metamorphic,khiabani2026agentic}.

Pendekatan \en{rules as code} memisahkan aturan dari kode dengan menuliskan peraturan sebagai spesifikasi yang dapat dieksekusi~\cite{merigoux2021catala,mowbray2023scaling,kowalski2022logicalenglish}. Representasi formal pengetahuan pajak juga telah dipakai untuk payroll dan akuntansi pajak~\cite{dumancic2021payroll,makowski2025dmn}, dan di Indonesia telah dibangun berbagai aplikasi perhitungan PPh 21~\cite{ramadhani2021grossup,yulistyawati2022payroll,kusuma2026pph21}. Namun, sejauh penelusuran penulis, belum ada pendekatan untuk PPh 21 yang sekaligus (a) memisahkan lapisan regulasi dari lapisan kebijakan perusahaan dengan resolusi konflik berbasis sifat aturan, (b) memvalidasi kedua rezim TER dan pra-TER dalam satu \en{knowledge base} berversi waktu, dan (c) menetapkan kebenaran sampai rupiah tanpa \en{domain expert}. Di sisi lain, \en{large language model} (LLM) mulai dipakai untuk penalaran hukum dan pajak, tetapi keandalannya masih terbatas, baik dalam penalaran undang-undang pajak~\cite{blairstanek2023gpt3,nay2024taxattorneys} maupun dalam menjawab pertanyaan faktual tentang hukum~\cite{dahl2024legalfictions}. Karena itu, penelitian ini memakai \en{inference engine} deterministik untuk menghitung, dan LLM hanya dipakai di luar \en{engine} untuk menyusun rancangan aturan yang wajib disetujui manusia.

Penelitian ini mengusulkan \en{knowledge-based system} (KBS) untuk perhitungan PPh Pasal 21 pegawai tetap tahun pajak 2023--2026. Kontribusi yang diusulkan adalah sebagai berikut.
\begin{enumerate}[leftmargin=*,nosep]
  \item \textbf{\en{Knowledge base} dua lapis dengan resolusi konflik berbasis asas hukum.} Hubungan antara regulasi dan kebijakan pemberi kerja dimodelkan eksplisit melalui sifat aturan (\texttt{wajib}, \texttt{default}, \texttt{opsional}, \texttt{tafsir}) dan asas \en{lex superior}, \en{lex specialis}, dan \en{lex posterior}. Dengan begitu, pelanggaran aturan wajib terdeteksi otomatis.
  \item \textbf{Versi waktu bitemporal.} \en{Valid time} memungkinkan satu \en{engine} menghitung rezim PER-16 dan rezim TER. \en{Knowledge time} memungkinkan pertanyaan ``berapa yang dihitung pada saat itu'' untuk regulasi yang berlaku surut. Adaptabilitas terhadap perubahan diukur dengan \en{impact precision}, \en{recall}, dan F1.
  \item \textbf{Perlakuan deklaratif atas \en{gross-up} sebagai \en{fixed point}.} Keberadaan solusi dijamin teorema Knaster--Tarski, solusi ganda dideteksi lewat \en{least fixed point} dan \en{greatest fixed point}, dan pilihan antar-solusi dibuat deterministik.
  \item \textbf{Metodologi validasi tanpa \en{domain expert}.} Kebenaran ditetapkan lewat \en{oracle pyramid}: contoh resmi dengan protokol erratum, \en{differential testing} terhadap \en{baseline} yang dibekukan lebih dulu, \en{metamorphic} dan \en{property-based testing}, serta \en{mutation testing}. Semuanya dijalankan di bawah kontrak presisi tanpa \en{floating point}.
  \item \textbf{Artefak terbuka.} Artefak yang dirilis meliputi \en{knowledge base} PPh 21 2023--2026, \en{test case} resmi terstruktur beserta katalog erratum, kalkulator pembanding tanpa KB, demo Streamlit, dan aplikasi web Laravel untuk admin \en{finance}.
\end{enumerate}

Pembahasan selanjutnya disusun sebagai berikut. Bagian~\ref{sec:tinjauan} membahas penelitian terkait. Bagian~\ref{sec:metodologi} menguraikan metodologi penelitian: domain, data, representasi \en{knowledge base}, \en{inference engine}, kontrak presisi, \en{baseline}, antarmuka dan aplikasi web, strategi validasi, serta desain eksperimen. Hasil eksperimen akan dilaporkan pada versi berikutnya.
